\documentclass[11pt,a4paper]{article}
\usepackage[utf8]{inputenc}
\usepackage[T1]{fontenc}
\usepackage[spanish,es-tabla,es-noquoting]{babel}
\usepackage[margin=2.2cm,top=2.3cm,bottom=2.3cm]{geometry}
\usepackage{xcolor}
\usepackage{graphicx}
\usepackage{booktabs}
\usepackage{tabularx}
\usepackage{array}
\usepackage{tcolorbox}
\usepackage{hyperref}
\usepackage{tikz}
\usetikzlibrary{positioning, arrows.meta}

% =========================================================================
% PALETA CORPORATIVA DE LUJO (AGRO-PREMIUM)
% =========================================================================
\definecolor{AgroGreen}{RGB}{11, 79, 38}        % Verde Esmeralda Agro
\definecolor{AgroGreenDark}{RGB}{6, 44, 21}     % Verde Bosque Profundo
\definecolor{AgroGold}{RGB}{217, 119, 6}        % Ámbar Cuero Vacuno
\definecolor{AgroDark}{RGB}{24, 24, 27}         % Carbón Ejecución
\definecolor{AgroCardBg}{RGB}{250, 250, 249}    % Fondo Suave
\definecolor{AgroBorder}{RGB}{229, 231, 235}    % Borde Fino

% Configuración de Hipervínculos
\hypersetup{
    colorlinks=true,
    linkcolor=AgroGreen,
    filecolor=AgroGold,
    urlcolor=AgroGreen,
    pdftitle={Dossier Comercial y Manual Operativo - BRK Agro Bolivia},
    pdfauthor={BRK Agro Tech Team},
    pdfsubject={Presentación Ejecutiva y Manual de Operaciones},
    pdfkeywords={BRK Agro, E-Commerce, ERP, Bolivia, UV50, POS, Ventas}
}

% Estilo de Cajas con tcolorbox
\newtcolorbox{propbox}[1][]{
  colback=white,
  colframe=AgroGreen,
  fonttitle=\bfseries\color{white},
  coltitle=white,
  colbacktitle=AgroGreen,
  arc=2.5mm,
  boxrule=1.2pt,
  left=4mm, right=4mm, top=3mm, bottom=3mm,
  title=#1
}

\newtcolorbox{rolebox}[2][]{
  colback=AgroCardBg,
  colframe=AgroGreenDark,
  fonttitle=\bfseries\color{white},
  colbacktitle=AgroGreenDark,
  arc=2mm,
  boxrule=1pt,
  left=3.5mm, right=3.5mm, top=2.5mm, bottom=2.5mm,
  title={#2},
  #1
}

\newtcolorbox{highlightbox}[1][]{
  colback=orange!5!white,
  colframe=AgroGold,
  fonttitle=\bfseries\color{AgroGold!90!black},
  colbacktitle=orange!15!white,
  arc=2mm,
  boxrule=0.9pt,
  left=3.5mm, right=3.5mm, top=2.5mm, bottom=2.5mm,
  title={Propuesta de Valor Estratégica},
  #1
}

% Estilos para Diagramas TikZ
\tikzset{
  flowstep/.style={
    rectangle, 
    rounded corners=4pt, 
    minimum width=6cm, 
    minimum height=0.95cm,
    align=center, 
    draw=AgroGreen, 
    fill=white, 
    line width=1.1pt,
    font=\small\bfseries\color{AgroDark}
  },
  flowarrow/.style={
    -{Stealth[scale=1.1]}, 
    line width=1.2pt, 
    draw=AgroGreen
  }
}

\begin{document}

% =========================================================================
% PORTADA EJECUTIVA DE PRESENTACIÓN COMERCIAL
% =========================================================================
\begin{titlepage}
    \centering
    \vspace*{0.5cm}
    
    {\Huge\textbf{\color{AgroGreenDark}BRK AGRO BOLIVIA}}\\[0.35cm]
    {\Large\textbf{\color{AgroGold}DOSSIER COMERCIAL, MANUAL OPERATIVO \& ROLES}}\\[0.6cm]
    
    \rule{\linewidth}{1.5pt}\\[0.4cm]
    {\LARGE\textbf{Ecosistema Digital E-Commerce \& ERP Agropecuario}}\\[0.25cm]
    {\large Plataforma de Alta Conversión en WhatsApp, Punto de Venta y Control Financiero}\\[0.4cm]
    \rule{\linewidth}{1.5pt}\\[1.2cm]
    
    \begin{propbox}[Ficha Ejecutiva de Implementación]
        \begin{itemize}
            \item \textbf{Plataforma en Producción:} \url{https://brk-agro-iota.vercel.app/}
            \item \textbf{Consola de Administración Central:} \url{https://brk-agro-iota.vercel.app/admin}
            \item \textbf{Infraestructura de Datos:} Servidor de Base de Datos en la Nube con Respaldo 24/7
            \item \textbf{Red de Entrega y Velocidad:} Red Global Edge con Carga Rápida en Celulares
            \item \textbf{Canal de Cierre Comercial:} Integración Nativa con WhatsApp Business
            \item \textbf{Versión del Sistema:} Entrega Oficial v1.0.0
        \end{itemize}
    \end{propbox}

    \vfill
    
    \begin{highlightbox}
    \textbf{Objetivo del Sistema:} Transformar las consultas en ventas cerradas, erradicar devoluciones por talla equivocada mediante algoritmos biométricos y dar al propietario el control exacto de su ganancia neta en tiempo real.
    \end{highlightbox}

    \vspace{0.8cm}
    {\normalsize\textbf{Elaborado para:} Dirección General, Socios Estratégicos y Equipo Comercial}\\[0.2cm]
    {\small Documento Confidencial de Presentación y Operaciones --- Bolivia, 2026}
\end{titlepage}

% =========================================================================
% ÍNDICE GENERAL
% =========================================================================
\tableofcontents
\newpage

% =========================================================================
% SECCIÓN 1: RESUMEN EJECUTIVO (POR QUÉ ESTA PLATAFORMA)
% =========================================================================
\section{Resumen Ejecutivo: La Propuesta de Valor}

El comercio de indumentaria y calzado agropecuario en Bolivia (camisas técnicas con protección solar UV50+, botinas de trabajo en cuero legítimo, jeans reforzados y accesorios de faena) enfrenta tres barreras críticas que frenan el crecimiento de los negocios tradicionales:

\begin{enumerate}
    \item \textbf{Fricción en la compra online:} El cliente del campo y ganadero no compra mediante pasarelas con tarjetas complejas; compra chateando por WhatsApp. La tienda web de BRK Agro estructura el pedido en segundos y lo envía con código único directo al chat del negocio.
    \item \textbf{Costo oculto de las devoluciones por talla:} Errar la talla significa fletes de encomienda pagados dos veces, clientes molestos y prendas estancadas. La \textit{Calculadora Biométrica de Tallas} elimina este problema sugiriendo la medida exacta.
    \item \textbf{Desconocimiento de la utilidad real:} Vender mucho no siempre significa ganar bien. El ERP desglosa las ventas brutas, descuenta el Costo de Mercadería Vendida (CMV) y los gastos operativos, mostrando la \textbf{Ganancia Neta Real} del propietario.
\end{enumerate}

% =========================================================================
% SECCIÓN 2: DIAGRAMAS DE FLUJO OPERATIVO (CÓMO FUNCIONA)
% =========================================================================
\section{Diagramas de Arquitectura y Flujo del Sistema}

A continuación se presentan los diagramas visuales que ilustran el recorrido de los clientes y la administración de los pedidos:

\subsection{Flujo 1: De la Visita Web al Cierre en WhatsApp}
Describe cómo un comprador selecciona productos, calcula su talla, activa el descuento por mayor y despacha su orden sin trabas:

\begin{center}
\begin{tikzpicture}
    \node (step1) [flowstep] {1. Navegación en Tienda \\ \scriptsize Catálogo UV50+ y Cuero Legítimo};
    \node (step2) [flowstep, below=0.8cm of step1] {2. Calculador de Talla \\ \scriptsize Estatura, peso y cm de pie};
    \node (step3) [flowstep, below=0.8cm of step2] {3. Descuento Mayorista \\ \scriptsize Automático desde 3 unidades};
    \node (step4) [flowstep, below=0.8cm of step3] {4. Checkout y Código \\ \scriptsize Generación de \#BRK-XXXX};
    \node (step5) [flowstep, below=0.8cm of step4, fill=AgroGreen!10] {5. Cierre en WhatsApp \\ \scriptsize Mensaje listo para empaque};

    \draw [flowarrow] (step1) -- (step2);
    \draw [flowarrow] (step2) -- (step3);
    \draw [flowarrow] (step3) -- (step4);
    \draw [flowarrow] (step4) -- (step5);
\end{tikzpicture}
\end{center}

\vspace{0.4cm}

\subsection{Flujo 2: Operación en Tienda Física (POS Mostrador)}
Describe cómo el personal cobra en el mostrador físico sin descuadrar inventario ni dinero:

\begin{center}
\begin{tikzpicture}
    \node (p1) [flowstep] {1. Selección Táctil \\ \scriptsize Prenda y Talla en Mostrador};
    \node (p2) [flowstep, below=0.8cm of p1] {2. Cobro Bimodal \\ \scriptsize Efectivo o Simple QR Bancario};
    \node (p3) [flowstep, below=0.8cm of p2] {3. Asiento de Kardex \\ \scriptsize Descuento de stock en talla};
    \node (p4) [flowstep, below=0.8cm of p3, fill=AgroGreen!10] {4. Emisión de Recibo \\ \scriptsize Entrega y registro en Caja};

    \draw [flowarrow] (p1) -- (p2);
    \draw [flowarrow] (p2) -- (p3);
    \draw [flowarrow] (p3) -- (p4);
\end{tikzpicture}
\end{center}

\newpage

\subsection{Flujo 3: Asistente de Sustitución Express}
El procedimiento técnico exclusivo para sustitución de prendas en sucursal:

\begin{center}
\begin{tikzpicture}
    \node (c1) [flowstep, draw=AgroGold] {Cliente solicita cambio \\ \scriptsize Presenta prenda en tienda};
    \node (c2) [flowstep, below=0.8cm of c1] {Reingreso al Inventario \\ \scriptsize Talla devuelta (+1 al stock)};
    \node (c3) [flowstep, below=0.8cm of c2] {Despacho de Nueva Talla \\ \scriptsize Talla entregada (-1 al stock)};
    \node (c4) [flowstep, below=0.8cm of c3, fill=green!10, draw=AgroGreen] {Cero Impacto en Caja \\ \scriptsize Registro Kardex sin descuadre};

    \draw [flowarrow] (c1) -- (c2);
    \draw [flowarrow] (c2) -- (c3);
    \draw [flowarrow] (c3) -- (c4);
\end{tikzpicture}
\end{center}

% =========================================================================
% SECCIÓN 3: MATRIZ DE ROLES Y ACCESOS DEL PERSONAL
% =========================================================================
\section{Matriz de Roles: Cómo Opera el Equipo el Sistema}

La plataforma está diseñada para que cada miembro del equipo tenga claridad sobre sus funciones cotidianas:

\begin{rolebox}{Rol 1: Propietario / Gerente General (Control Total)}
\textbf{Acceso:} Panel ERP completo desde el conmutador inferior o \texttt{/admin}.\\
\textbf{Responsabilidades Diarias:}
\begin{itemize}
    \item \textbf{Supervisión Financiera:} Lectura del Dashboard (Ventas Brutas vs Costo CMV vs Utilidad Neta Real).
    \item \textbf{Aprobación de Compras:} Visualización de la matriz de tallas en semáforo para autorizar compras a fábrica en tallas agotadas.
    \item \textbf{Arqueo y Cierre:} Validación cruzada entre efectivo en caja y depósitos bancarios QR.
    \item \textbf{Control de Marca:} Modificación del teléfono de WhatsApp comercial, logotipo y anuncios.
\end{itemize}
\end{rolebox}

\begin{rolebox}{Rol 2: Vendedor de Mostrador / Cajero (Punto Físico)}
\textbf{Acceso:} Pestaña \textit{Pedidos y POS} en el panel administrativo.\\
\textbf{Responsabilidades Diarias:}
\begin{itemize}
    \item Cobro rápido de prendas seleccionando talla y medio de pago (Efectivo / QR).
    \item Emisión de comprobantes de venta impresos para el comprador presencial.
    \item Asentamiento obligatorio de cualquier gasto menor de caja chica (fletes, cintas, refrigerios).
    \item Entrega del reporte de caja al culminar la jornada.
\end{itemize}
\end{rolebox}

\begin{rolebox}{Rol 3: Encargado de Logística y Envíos (Flotas a Provincias)}
\textbf{Acceso:} Pestaña \textit{Centro de Impresiones} y bandeja de pedidos.\\
\textbf{Responsabilidades Diarias:}
\begin{itemize}
    \item Recepción de pedidos originados en la web con código correlativo (\texttt{\#BRK-XXXX}).
    \item Verificación de stock y aprobación en 1 clic para descontar la mercadería.
    \item Impresión del \textbf{Rótulo Gigante de Encomienda} para despacho en empresas de transporte terrestre y flotas interprovinciales.
\end{itemize}
\end{rolebox}

\begin{rolebox}{Rol 4: Comprador Final y Clientes Mayoristas}
\textbf{Acceso:} Tienda pública en línea (\url{https://brk-agro-iota.vercel.app/}).\\
\textbf{Experiencia del Usuario:}
\begin{itemize}
    \item Explora el catálogo con fichas técnicas de protección UV50+ y calzado de cuero legítimo.
    \item Usa la calculadora para conocer su talla exacta de camisa o botina.
    \item Obtiene descuento automático si adquiere 3 o más prendas.
    \item Cierra el trato directamente con la empresa a través de WhatsApp.
\end{itemize}
\end{rolebox}

\newpage

% =========================================================================
% SECCIÓN 4: CREDENCIALES Y ACCESOS OFICIALES DEL SISTEMA
% =========================================================================
\section{Credenciales de Inicio de Sesión y Accesos Oficiales}

Para que la dirección y el personal puedan evaluar y operar el sistema de inmediato, se han habilitado las credenciales maestras y los perfiles de acceso correspondientes:

\subsection{Tabla Maestra de Credenciales por Rol}

\begin{table}[h!]
\centering
\small
\begin{tabularx}{\linewidth}{p{3.5cm} p{3.8cm} p{2.6cm} X}
\toprule
\textbf{Rol / Cargo} & \textbf{Usuario / Correo} & \textbf{Clave / PIN} & \textbf{Permisos y Pestaña Principal} \\
\midrule
\textbf{Gerente General / Dueño} & \texttt{gerencia@brkagro.bo} & \texttt{BRK2026*Admin} & Acceso total a todas las pestañas: Finanzas, Ventas, Kardex, Precios y Ajustes de Marca. \\
\addlinespace
\textbf{Cajero / Venta Mostrador} & \texttt{caja@brkagro.bo} & \texttt{CajaBRK2026} \newline \textit{PIN: 2026} & Operación exclusiva de Punto de Venta (POS), cobro bimodal (Efectivo/QR) y arqueo. \\
\addlinespace
\textbf{Encargado de Logística} & \texttt{despacho@brkagro.bo} & \texttt{EnviosBRK2026} \newline \textit{PIN: 3030} & Aprobación de pedidos web y generación del Rótulo Gigante para encomiendas interprovinciales. \\
\addlinespace
\textbf{Bodega / Inventario} & \texttt{almacen@brkagro.bo} & \texttt{StockBRK2026} \newline \textit{PIN: 4040} & Matriz de stock por tallas (S--XXL / 38--44), recepción de fábrica y kardex físico. \\
\addlinespace
\textbf{Cliente Final / Mayorista} & \textit{Acceso Público Libre} & \textit{Sin clave} & Catálogo en línea, calculadora biométrica de tallas y checkout directo a WhatsApp. \\
\bottomrule
\end{tabularx}
\caption{Matriz oficial de credenciales para demostración y puesta en marcha.}
\end{table}

\subsection{Instrucciones para Iniciar Sesión en la Plataforma}

El acceso a la plataforma está optimizado para facilitar la demostración ejecutiva con el cliente:

\begin{propbox}[Métodos de Ingreso al Sistema]
\begin{enumerate}
    \item \textbf{Ingreso Directo a la Tienda Web (Cliente):}
    \begin{itemize}
        \item Ingrese a: \url{https://brk-agro-iota.vercel.app/}
        \item Navegue por el catálogo, seleccione prendas y pruebe la calculadora de tallas.
    \end{itemize}
    
    \item \textbf{Ingreso a la Consola Administrativa ERP (Gerencia y Personal):}
    \begin{itemize}
        \item Ingrese directamente a: \url{https://brk-agro-iota.vercel.app/admin}
        \item O utilice el \textbf{Botón Conmutador Flotante} situado en la esquina inferior izquierda de la pantalla, que permite alternar instantáneamente entre la Tienda y el ERP con un solo toque.
    \end{itemize}

    \item \textbf{Selector de Identidad Visual (Temas):}
    \begin{itemize}
        \item En la barra superior o conmutador, puede alternar entre el estilo contemporáneo (\textit{BRK Brasil Dark}) y el estilo tradicional de estancia (\textit{BRK Bolivia Country}).
    \end{itemize}
\end{enumerate}
\end{propbox}

\subsection{Protocolo de Registro y Alta de Nuevos Empleados}

El sistema cuenta con un módulo de \textbf{Gestión de Personal y Control de Accesos} que permite al propietario dar de alta, suspender o gestionar a su equipo de trabajo sin requerir conocimientos técnicos:

\begin{center}
\begin{tikzpicture}
    \node (e1) [flowstep] {1. Panel de Ajustes \\ \scriptsize Pestaña Marca \& Sistema};
    \node (e2) [flowstep, below=0.7cm of e1] {2. Formulario de Alta \\ \scriptsize Botón «+ Registrar Empleado»};
    \node (e3) [flowstep, below=0.7cm of e2] {3. Asignación de Rol \& PIN \\ \scriptsize Permisos (POS, Despacho, Bodega)};
    \node (e4) [flowstep, below=0.7cm of e3, fill=AgroGreen!10] {4. Habilitación Inmediata \\ \scriptsize Firma de ventas y arqueos de caja};

    \draw [flowarrow] (e1) -- (e2);
    \draw [flowarrow] (e2) -- (e3);
    \draw [flowarrow] (e3) -- (e4);
\end{tikzpicture}
\end{center}

\vspace{0.3cm}

\begin{propbox}[Procedimiento Paso a Paso para Registrar un Nuevo Usuario]
\begin{enumerate}
    \item \textbf{Ubicación del Módulo:}
    \begin{itemize}
        \item Ingrese a la consola de administración en \url{https://brk-agro-iota.vercel.app/admin}.
        \item Seleccione la pestaña \textbf{«Marca \& Sistema»} y desplácese hasta la sección \textbf{«Equipo de Trabajo \& Control de Accesos (Personal)»}.
    \end{itemize}

    \item \textbf{Apertura del Registro:}
    \begin{itemize}
        \item Presione el botón superior derecho \textbf{\texttt{[+ Registrar Empleado]}}.
        \item Se desplegará la ventana modal de alta de personal.
    \end{itemize}

    \item \textbf{Diligenciamiento de Datos Requeridos:}
    \begin{itemize}
        \item \textbf{Nombre Completo:} Nombre y apellidos del colaborador (ej. \textit{Juan Carlos Mamani}).
        \item \textbf{Cargo / Rol Operativo:} Seleccione entre:
        \begin{itemize}
            \item \textit{Cajero / POS:} Habilita cobros en mostrador físico, arqueo de caja y cambio de tallas.
            \item \textit{Logística / Despacho:} Habilita aprobación de pedidos WhatsApp e impresión de rótulos de flota.
            \item \textit{Bodega / Inventario:} Habilita control matricial de tallas y recepción de fábrica.
            \item \textit{Gerente General:} Otorga control irrestricto sobre finanzas, precios y configuración.
        \end{itemize}
        \item \textbf{PIN de Seguridad (4 dígitos):} Código numérico confidencial (ej. \texttt{2026}) con el que el colaborador valida sus operaciones.
        \item \textbf{Celular / WhatsApp:} Número de contacto para coordinación operativa.
        \item \textbf{Correo Electrónico (Opcional):} Dirección de correo corporativo o personal.
    \end{itemize}

    \item \textbf{Guardado y Confirmación:}
    \begin{itemize}
        \item Presione \textbf{«Guardar Empleado»}. El sistema le asigna un identificador único (\texttt{emp-XXXX}) y lo incorpora inmediatamente a la matriz de personal.
    \end{itemize}

    \item \textbf{Gestión de Estados (Activar / Suspender / Eliminar):}
    \begin{itemize}
        \item \textbf{Suspensión Temporal:} El ícono de verificación permite alternar entre estado \textit{Activo} e \textit{Inactivo} sin borrar el historial contable del empleado.
        \item \textbf{Baja Definitiva:} El ícono de papelera permite eliminar el perfil cuando un colaborador ya no forma parte de la empresa.
    \end{itemize}
\end{enumerate}
\end{propbox}

\begin{highlightbox}
\textbf{Trazabilidad y Auditoría Contable:} Cada vez que un cajero realiza una venta en mostrador, efectúa un arqueo de caja chica o emite un cambio de talla express, el sistema asienta su nombre y rol en los registros de auditoría (\texttt{sellerName}, \texttt{closedBy}), garantizando que el propietario sepa con exactitud quién procesó cada boliviano en la tienda.
\end{highlightbox}

\subsection{Infraestructura de Datos y Servidor Cloud}
El sistema se encuentra interconectado a su clúster de base de datos corporativo en la nube:
\begin{itemize}
    \item \textbf{Servidor Principal:} Supabase Cloud Enterprise (\texttt{Brk\_Agrobd})
    \item \textbf{Identificador de Proyecto:} \texttt{rwnuknrqjolwvvgfandg}
    \item \textbf{Esquema de Tablas:} 8 tablas sincronizadas en tiempo real (\texttt{products}, \texttt{product\_variants}, \texttt{whatsapp\_orders}, \texttt{sales}, \texttt{cash\_movements}, \texttt{kardex\_movements}, \texttt{size\_exchanges}, \texttt{store\_settings}).
    \item \textbf{Respaldo y Seguridad:} Respaldos automáticos continuos con cifrado SSL en reposo y tránsito.
\end{itemize}

\begin{highlightbox}
\textbf{Resiliencia Operativa en Campo:} En caso de cortes imprevistos de conexión a internet en la estancia o sucursal provincial, el módulo de ventas y caja cuenta con persistencia local en caché; el cajero continúa cobrando sin interrupción y los registros se sincronizan automáticamente al restablecerse la red.
\end{highlightbox}

% =========================================================================
% SECCIÓN 5: CAPACIDAD DE ADAPTACIÓN Y PERSONALIZACIÓN TOTAL
% =========================================================================
\section{Capacidad de Adaptación: Todo lo que Puede Cambiarse}

Uno de los mayores atractivos para la empresa es que el sistema es completamente \textbf{White-Label y modular}, lo que significa que el negocio puede cambiar sus parámetros sin depender de soporte técnico continuo:

\subsection{¿Qué puede modificar el dueño en tiempo real?}
Desde la pestaña \textbf{Ajustes / Marca} del panel administrativo, el propietario puede cambiar en 1 clic:
\begin{itemize}
    \item \textbf{Número de WhatsApp de Ventas:} Si cambia de vendedor o de línea telefónica, coloca el nuevo número y todos los pedidos web se redirigen inmediatamente.
    \item \textbf{Nombre Comercial y Eslogan:} Adaptable para campañas estacionales (ej. \textit{«Feria Ganadera FEXPOCRUZ»}).
    \item \textbf{Barra de Anuncios Superior:} Para publicar promociones temporales (ej. \textit{«Envíos gratis por compras mayores a Bs. 500»}).
    \item \textbf{Cuentas Bancarias y QR Oficial:} Actualización de números de cuenta para transferencias directas.
    \item \textbf{Precios de Venta y Costos:} Edición inmediata del precio al público, precio mayorista y costo de fábrica para mantener la rentabilidad actualizada ante variaciones de mercado.
    \item \textbf{Catálogo de Productos:} Alta y baja de productos, fotografías, tallas y descripciones con guardado instantáneo.
\end{itemize}

\subsection{Escalabilidad Futura}
El sistema está construido con estándares modernos que permiten integrar en etapas posteriores:
\begin{itemize}
    \item Múltiples sucursales físicas (Santa Cruz, Cochabamba, Beni, Tarija).
    \item Facturación electrónica computarizada en línea (SIAT).
    \item Lectores de código de barras para pistoleros de inventario en bodega.
\end{itemize}

\newpage

% =========================================================================
% SECCIÓN 6: MATRIZ DE REGLAS OPERATIVAS (QUÉ SÍ VS QUÉ NO)
% =========================================================================
\section{Matriz de Reglas Operativas: Buenas Prácticas}

Para garantizar que el sistema mantenga su precisión y la contabilidad esté blindada:

\begin{table}[h!]
\centering
\small
\begin{tabularx}{\linewidth}{p{3.2cm} X X}
\toprule
\textbf{Aspecto} & \textbf{[SI] Qué SÍ debe hacerse} & \textbf{[NO] Qué NO debe hacerse} \\
\midrule
\textbf{Gestión de Tallas} & Consultar estatura y peso del cliente para confirmar la talla sugerida por la calculadora. & Despachar tallas adivinando sin corroborar la complexión del comprador. \\
\addlinespace
\textbf{Cambios en Tienda} & Utilizar siempre el botón \textit{Cambio de Talla Express} al recibir una prenda devuelta. & Entregar otra prenda de palabra sin registrar la salida del stock. \\
\addlinespace
\textbf{Ventas de Mostrador} & Registrar cada venta en el POS señalando si fue en Efectivo o Simple QR. & Guardar dinero en cajón sin emitir recibo o nota en el sistema. \\
\addlinespace
\textbf{Caja Chica} & Asentar en el módulo de caja todo egreso por pequeño que sea (bolsas, cinta, pasajes). & Retirar efectivo para gastos del personal sin registrar el asiento. \\
\addlinespace
\textbf{Parámetros de Marca} & Cambiar teléfonos y logotipos desde la pestaña \textit{Ajustes/Marca}. & Intentar alterar archivos internos o base de datos a mano. \\
\bottomrule
\end{tabularx}
\end{table}

% =========================================================================
% SECCIÓN 7: GUION MAESTRO PARA CONVENCER AL CLIENTE (5 MINUTOS)
% =========================================================================
\section{Guion Maestro de Cierre Comercial (Para el Cliente)}

Utilice este guion estructurado al momento de presentar la plataforma para demostrar su valor indiscutible:

\begin{tcolorbox}[colback=white,colframe=AgroGreen!80!black,title=\textbf{Paso 1: Demostrar Presencia y Nivel Gráfico (Minuto 1)}]
\textit{«Mire la plataforma en vivo en \texttt{brk-agro-iota.vercel.app}. Esta tienda no parece una plantilla genérica; tiene la presencia visual y la categoría de las grandes marcas de indumentaria agropecuaria de Brasil (BRK Brasil). Está concebida para que el productor ganadero o agrónomo sienta confianza inmediata desde su celular.»}
\end{tcolorbox}

\begin{tcolorbox}[colback=white,colframe=AgroGreen!80!black,title=\textbf{Paso 2: La Solución Definitiva a las Devoluciones (Minuto 2)}]
\textit{«El principal temor de quien compra ropa o botinas por internet es: '¿Y si me queda chica?'. Mire esto: el cliente ingresa su estatura y su peso, y el sistema le indica con exactitud matemática qué talla debe pedir. Esto le ahorra a usted hasta un 70\% de gastos de encomienda por cambios y reclamos.»}
\end{tcolorbox}

\begin{tcolorbox}[colback=white,colframe=AgroGreen!80!black,title=\textbf{Paso 3: Venta por Mayor y Cierre en WhatsApp (Minuto 3)}]
\textit{«Si una estancia o empresa compra 3 o más unidades, el sistema activa automáticamente el precio por mayor y le muestra al cliente cuánto está ahorrando en el acto. Y al tocar finalizar pedido, no le pide tarjetas difíciles: le arma el pedido perfecto con código correlativo (\#BRK-XXXX) directo a su WhatsApp.»}
\end{tcolorbox}

\begin{tcolorbox}[colback=white,colframe=AgroGreen!80!black,title=\textbf{Paso 4: El Control Financiero y Logístico en el ERP (Minuto 4)}]
\textit{«En el panel administrativo usted ve en una sola pantalla cuánto vendió, cuánto costó la mercadería y su **utilidad neta real limpia**. Tiene un POS para ventas en mostrador, control de inventario por tallas, cambio de tallas express y rótulos gigantes para despachar en las flotas a provincia sin extravíos.»}
\end{tcolorbox}

\begin{tcolorbox}[colback=white,colframe=AgroGreen!80!black,title=\textbf{Paso 5: Cierre de Negocio y Entrega Inmediata (Minuto 5)}]
\textit{«La plataforma está lista, probada y funcionando en la nube hoy mismo. Con este sistema su negocio deja de operar en cuadernos o memorias y se convierte en una empresa profesional y rentable.»}
\end{tcolorbox}

\vspace{0.8cm}
\begin{center}
\rule{0.6\linewidth}{0.6pt}\\[0.2cm]
\small\textbf{BRK Agro Bolivia} \textbar{} Ecosistema Comercial E-Commerce \& ERP\\
Cochabamba \& Santa Cruz, Bolivia --- Octubre 2026
\end{center}

\end{document}
